\chapter{Estado actual del proyecto}

\section{Situación de arranque}

El proyecto ya dispone de estructura, automatización básica y dos ediciones sincronizadas.
La siguiente puerta de calidad exige un corpus real de PDF para poder auditar tipologías de
ejercicios, detectar duplicados y construir la cobertura trazable.

\begin{theorybox}{Condición necesaria para continuar}
Sin fuentes no puede garantizarse la correspondencia entre los ejercicios originales y las
secciones del cuaderno. Por ese motivo, la taxonomía presente en \texttt{data/exercise\_taxonomy.yml}
se considera solo una hipótesis de trabajo.
\end{theorybox}

\begin{methodbox}{Qué ocurrirá cuando lleguen los PDF}
\begin{enumerate}
  \item Se ejecutará la auditoría con hashes SHA-256 y huella textual normalizada.
  \item Se clasificará cada archivo por rol, nivel y fiabilidad de extracción.
  \item Se construirá la matriz de cobertura \texttt{fuente -> tipo -> sección}.
  \item Se aprobará un capítulo piloto antes de redactar el resto del libro.
\end{enumerate}
\end{methodbox}

\begin{solvedexample}{Comprobación de visibilidad entre ediciones}
En la edición del alumnado deben aparecer la explicación general y las respuestas breves,
pero no el desarrollo completo de la solución. En la edición del profesorado sí debe verse
la solución completa. Esta sección actúa como prueba de esa separación editorial.
\end{solvedexample}

\begin{guidedexercise}{Identificar el bloqueo actual}
Lee el estado del repositorio y contesta:
\begin{enumerate}
  \item ¿Existe algún PDF fuente en la carpeta?
  \item ¿Puede superarse la Fase 1 sin ese corpus?
\end{enumerate}
\end{guidedexercise}

\begin{shortanswer}{Ejercicio guiado}
\begin{enumerate}
  \item No, actualmente hay cero PDF localizados.
  \item No, porque la auditoría quedaría vacía y sin tipologías trazables.
\end{enumerate}
\end{shortanswer}

\begin{fullsolution}{Ejercicio guiado}
La búsqueda recursiva del proyecto no ha encontrado ningún archivo con extensión \texttt{.pdf}.
Como la Fase 1 exige inventario, clasificación y detección de duplicados del corpus, la
puerta de calidad no puede darse por superada mientras falten las fuentes originales.
\end{fullsolution}

\begin{commonerror}{Confundir plantilla con libro terminado}
Un error razonable sería pensar que, porque ya existen archivos \texttt{tex/}, \texttt{docs/} y \texttt{data/},
el cuaderno puede darse por hecho. No es así: la estructura técnica no sustituye la auditoría
real del material de partida.
\end{commonerror}

\begin{challenge}{Siguiente acción útil}
Incorpora los PDF originales a la carpeta raíz o a \texttt{sources/} y vuelve a ejecutar la
auditoría. Solo entonces tendrá sentido ampliar la taxonomía y producir contenido didáctico
definitivo.
\end{challenge}

\begin{selfcheck}
\begin{itemize}
  \item Sé localizar el bloqueo actual del proyecto.
  \item Entiendo por qué no conviene inventar cobertura sin corpus.
  \item Puedo reanudar el flujo con \texttt{scripts/tasks.ps1 audit}.
\end{itemize}
\end{selfcheck}
